%==============================================================================
\section{madGLP-to-Networking Mapping}
\label{sec:mapping}
%==============================================================================
%------------------------------------------------------------------------------
\subsection{Agent Identity Mapping}
\label{sec:id-mapping}
%------------------------------------------------------------------------------
Each madGLP agent identifier (e.g., \code{alice}, \code{bob}) is mapped to an Ed25519 public key. The GLP runtime maintains a bidirectional mapping:
\begin{itemize}
\item \func{resolve(agentId)} $\to$ \code{PubKey}: given a madGLP agent name, return the public key.
\item \func{resolve(pubKey)} $\to$ \code{AgentId}: given a public key, return the madGLP agent name.
\end{itemize}
This mapping is established during peer discovery: when agent $p$ discovers agent $q$ via the networking layer, the discovery handshake provides $q$'s public key, and the runtime registers the association between $q$'s agent name and public key.
For cold-calls, the sender of the network payload must know the target agent's public key. This is obtained either through peer discovery (e.g., BLE proximity) or through out-of-band peer exchange (Section~\ref{sec:discovery}).  In an IP peer-link cold-call, $A$ creates the signed invite, but the holder~$B$ becomes the first network sender: the invite supplies $A_{\mathit{pk}}$, allowing~$B$ to address $A$'s index-0 serializer.
%------------------------------------------------------------------------------
\subsection{Message Serialization}
\label{sec:serialization}
%------------------------------------------------------------------------------
A madGLP message $(G := T^\uparrow, Q)$ must be serialized into a byte payload for transmission. The payload encodes two components: the global variable name $G$ and the globalized term $T^\uparrow$.
\subsubsection{Global Variable Name Encoding}
A global name $\_w(p, i)$ or $\_r(p, i)$ is encoded as:
\begin{longtable}{L{3cm} L{3cm} L{7cm}}
\toprule
\rowcolor{headerblue}
\textcolor{white}{\textbf{Field}} &
\textcolor{white}{\textbf{Size}} &
\textcolor{white}{\textbf{Description}} \\
\midrule
\endhead
\code{tag} & 1 byte & \code{0x01} for $\_w$, \code{0x02} for $\_r$ \\
\rowcolor{rowgray}
\code{agent} & 32 bytes & Ed25519 public key of agent $p$ \\
\code{index} & 4 bytes & Index $i$ (big-endian unsigned 32-bit integer) \\
\bottomrule
\end{longtable}
Total: 37 bytes per global name.
\subsubsection{Globalized Term Encoding}
A globalized term $T^\uparrow$ is a ground term in which all variables have been replaced by global names. It is a standard term (functors, atoms, integers, lists) that may contain embedded global names. The encoding uses a tagged format:
\begin{longtable}{L{3cm} L{2.5cm} L{7.5cm}}
\toprule
\rowcolor{headerblue}
\textcolor{white}{\textbf{Tag}} &
\textcolor{white}{\textbf{Value}} &
\textcolor{white}{\textbf{Encoding}} \\
\midrule
\endhead
Atom & \code{0x10} & Length-prefixed UTF-8 string \\
\rowcolor{rowgray}
Integer & \code{0x11} & 8-byte signed big-endian \\
Global name & \code{0x12} & 37-byte global name encoding (as above) \\
\rowcolor{rowgray}
Functor & \code{0x20} & Tag, length-prefixed name, arity (2 bytes), then each argument recursively \\
List cell & \code{0x21} & Tag, head (recursive), tail (recursive) \\
\rowcolor{rowgray}
Empty list & \code{0x22} & Tag only \\
\bottomrule
\end{longtable}
\subsubsection{Message Envelope}
The complete serialized message is:
\begin{longtable}{L{3cm} L{3cm} L{7cm}}
\toprule
\rowcolor{headerblue}
\textcolor{white}{\textbf{Field}} &
\textcolor{white}{\textbf{Size}} &
\textcolor{white}{\textbf{Description}} \\
\midrule
\endhead
\code{version} & 1 byte & Protocol version (currently \code{0x01}) \\
\rowcolor{rowgray}
\code{msg\_type} & 1 byte & \code{0x01} for regular link, \code{0x02} for cold-call (index-0 serializer) \\
\code{global\_name} & 37 bytes & Target global variable name $G$ \\
\rowcolor{rowgray}
\code{term\_length} & 4 bytes & Length of the serialized term in bytes \\
\code{term\_data} & variable & Serialized globalized term $T^\uparrow$ \\
\bottomrule
\end{longtable}
%------------------------------------------------------------------------------
\subsection{The Send Operation}
\label{sec:send-op}
%------------------------------------------------------------------------------
When the madGLP runtime executes the Send transaction---moving message $(G := T^\uparrow, Q)$ from agent $p$'s outgoing set $M_p$ to the communication channel---the networking layer performs:
\begin{enumerate}
\item \textbf{Resolve destination}: Look up the public key for agent $Q$ via \func{resolve(Q)}.
\item \textbf{Serialize}: Encode the message $(G, T^\uparrow)$ into a byte payload according to the format in Section~\ref{sec:serialization}.
\item \textbf{Transmit}: Call \func{send(pubKey\_of\_Q, payload)}.
\end{enumerate}
If the destination agent is not currently reachable, the message is queued locally until the peer becomes reachable. This preserves madGLP's fair message delivery assumption (Section~\ref{sec:assumptions}).
%------------------------------------------------------------------------------
\subsection{The Receive Operation}
\label{sec:receive-op}
%------------------------------------------------------------------------------
When a payload arrives, the networking layer delivers it to the madGLP runtime:
\begin{enumerate}
\item \textbf{Identify sender}: The \func{onReceive} callback provides the sender's public key. Resolve to a madGLP agent identifier via \func{resolve(pubKey)}.
\item \textbf{Deserialize}: Decode the byte payload to recover the global variable name $G$ and the globalized term $T^\uparrow$.
\item \textbf{Dispatch}: Pass $(G, T^\uparrow)$ to the madGLP Receive transaction, which performs localization and assignment as specified in the madGLP definition.
\end{enumerate}
